\booksection{第十一章附：专业影像技术知识}

\chapterintro{想要真正理解手机摄影和视频的奥秘，你需要了解一些专业的影像技术知识。这些知识不仅能帮你更好地使用手机相机，还能让你在选择设备时更有判断力。}

\booksubsection{一、分辨率、帧率与码率：视频质量的三要素}

\subsubsection*{分辨率（Resolution）：}

分辨率决定了画面的清晰度，用像素数表示。

\begin{itemize}
  \item \textbf{720P (HD)}：1280×720像素，入门级清晰度。
  \item \textbf{1080P (Full HD)}：1920×1080像素，目前最主流的分辨率。
  \item \textbf{2K (QHD)}：2560×1440像素，比1080P更清晰。
  \item \textbf{4K (UHD)}：3840×2160像素，超高清，细节丰富。
  \item \textbf{8K}：7680×4320像素，未来标准，目前设备支持有限。
\end{itemize}

分辨率越高，画面越清晰，但文件也越大，对存储和处理能力要求越高。

\subsubsection*{帧率（Frame Rate）：}

帧率决定了画面的流畅度，用每秒帧数（fps）表示。

\begin{itemize}
  \item \textbf{24fps}：电影标准帧率，有独特的"电影感"。
  \item \textbf{30fps}：电视和网络视频常用，流畅自然。
  \item \textbf{60fps}：高帧率，适合运动场景和游戏录制，非常流畅。
  \item \textbf{120fps/240fps}：超高帧率，可用于慢动作视频。
\end{itemize}

帧率越高，画面越流畅，慢动作效果越好，但文件也越大。

\subsubsection*{码率（Bitrate）：}

码率决定了视频的数据密度，用每秒数据量（Mbps）表示。

\begin{itemize}
  \item 码率越高，画质越好，细节保留越多，但文件也越大。
  \item 码率过低会导致画面出现马赛克、色块等压缩痕迹。
  \item 专业视频拍摄通常需要高码率以保留更多后期处理空间。
\end{itemize}

\booksubsection{二、视频编码格式：H.264、H.265与ProRes}

\subsubsection*{H.264 (AVC)：}

\begin{itemize}
  \item \textbf{特点}：最广泛使用的编码格式，兼容性极好。
  \item \textbf{优点}：几乎所有设备都能播放，压缩效率适中。
  \item \textbf{缺点}：相比新格式，压缩效率较低，文件较大。
  \item \textbf{适用场景}：网络视频分享、普通用户日常拍摄。
\end{itemize}

\subsubsection*{H.265 (HEVC)：}

\begin{itemize}
  \item \textbf{特点}：H.264的升级版，压缩效率提升约50\%。
  \item \textbf{优点}：相同画质下文件更小，支持更高分辨率和帧率。
  \item \textbf{缺点}：编码和解码需要更强的处理器，部分老设备不兼容。
  \item \textbf{适用场景}：4K视频拍摄、存储空间有限的场景。
\end{itemize}

\subsubsection*{Apple ProRes：}

\begin{itemize}
  \item \textbf{特点}：苹果开发的专业视频编码，几乎无损压缩。
  \item \textbf{优点}：画质极高，后期处理空间大，色彩准确。
  \item \textbf{缺点}：文件极大，需要大量存储空间，处理需要专业设备。
  \item \textbf{适用场景}：专业视频制作、电影级拍摄。
\end{itemize}

\subsubsection*{ProRes RAW：}

\begin{itemize}
  \item \textbf{特点}：结合了ProRes和RAW格式的优势。
  \item \textbf{RAW是什么}：未经处理的原始传感器数据，保留最大信息量。
  \item \textbf{优点}：后期调整空间极大，可以改变曝光、白平衡等参数而不损失画质。
  \item \textbf{缺点}：文件极其巨大，需要专业后期软件处理。
  \item \textbf{适用场景}：高端专业制作。iPhone 15 Pro系列及以后支持拍摄Apple ProRes视频（最高4K 60fps），iPhone 17 Pro进一步支持ProRes RAW直接机内录制。
\end{itemize}

\booksubsection{三、HDR：高动态范围技术}

HDR（High Dynamic Range）技术可以同时保留画面中最亮和最暗部分的细节。

\begin{itemize}
  \item \textbf{HDR10}：最基础的HDR标准，兼容性好。
  \item \textbf{HDR10+}：三星推动的标准，支持动态元数据。
  \item \textbf{Dolby Vision}：杜比开发的顶级HDR标准，画质最佳，但需要授权。
  \item \textbf{HLG}：广播级HDR标准，兼容普通显示器。
\end{itemize}

拍摄HDR视频时，需要支持HDR的显示器才能看到完整效果。

\booksubsection{四、色彩空间与色深}

\subsubsection*{色彩空间：}

\begin{itemize}
  \item \textbf{sRGB}：标准色彩空间，兼容性最好。
  \item \textbf{Display P3}：苹果推广的广色域标准，色彩更丰富。
  \item \textbf{Rec.2020}：HDR视频的标准色彩空间，色域最广。
\end{itemize}

\subsubsection*{色深：}

\begin{itemize}
  \item \textbf{8-bit}：1670万色，普通视频标准。
  \item \textbf{10-bit}：10.7亿色，专业视频标准，色彩过渡更平滑。
  \item \textbf{12-bit}：680亿色，高端专业标准。
\end{itemize}

\begin{tcbinfo}{实用建议}
\begin{itemize}
  \item 日常拍摄：1080P/4K + 30fps + H.265编码，平衡画质和文件大小。
  \item 运动拍摄：1080P/4K + 60fps，保证流畅度。
  \item 专业创作：4K + ProRes（如果支持），保留最大后期空间。
  \item 存储建议：高码率视频很占空间，建议定期备份到电脑或云盘。
\end{itemize}

\end{tcbinfo}

\begin{tcbpractice}{去探索一下：你的手机视频设置}
\begin{enumerate}
  \item 打开相机App，切换到视频模式。
  \item 查看分辨率和帧率选项，尝试不同的组合。
  \item 如果支持，查看是否有HDR、ProRes等高级选项。
  \item 拍摄几段不同设置的视频，对比画质和文件大小。
\end{enumerate}

\end{tcbpractice}

\booksubsection{五、RAW格式与HEIF/JPEG对比}

\subsubsection*{什么是RAW格式：}

RAW格式是相机传感器记录的“原始数据”——未经任何压缩和处理的纯原始信息。它就像生鲜食材，需要后期“烹饪”才能成品。手机因为存储和算力有限，默认拍出来的都是已经“烹饪好”的JPEG或HEIF，RAW需要专门开启。

\subsubsection*{RAW vs HEIF vs JPEG对比：}

\begin{itemize}
  \item \textbf{JPEG}：最老的格式，兼容性最好（任何设备都能打开）。8-bit色深，有损压缩，文件最小（一张照片2-5MB）。后期空间有限，强行提亮暗部容易出现色块。
  \item \textbf{HEIF/HEIC}：苹果主推的新格式，同样画质下文件只有JPEG的一半（一张照片1-3MB）。支持10-bit色深和透明通道。缺点是部分老设备和Windows需要安装插件才能打开。
  \item \textbf{RAW}：未经压缩的原始数据，文件最大（一张照片20-100MB）。色深通常12-16bit，后期空间最大。缺点是必须用专业软件（Lightroom、Snapseed等）才能打开和处理，无法直接分享。
\end{itemize}

\subsubsection*{为什么专业摄影师要拍RAW：}

\begin{itemize}
  \item \textbf{白平衡可后期无损调整}：JPEG拍错了白平衡（比如白炽灯下偏黄），后期纠正会损失画质；RAW可以重新设置色温，完全无损。
  \item \textbf{曝光容差大}：JPEG欠曝2档的暗部基本是死黑一片，提亮后全是噪点；RAW能从欠曝3-4档的暗部"救"出可用细节。过曝同理。
  \item \textbf{暗部和高光细节更丰富}：RAW的动态范围远大于JPEG，亮部不过曝、暗部不死黑的余地大得多。
  \item \textbf{锐化、降噪、镜头矫正都在后期可控}：相机的机内处理是"一锤子买卖"，RAW可以反复调整，每次都从原始数据出发。
\end{itemize}

\subsubsection*{手机如何开启RAW：}

\begin{itemize}
  \item \textbf{iPhone ProRAW}：仅iPhone 12 Pro及以上Pro机型支持。打开"设置 相机 格式"中开启"Apple ProRAW"。然后在相机界面顶部会出现一个"RAW"按钮，点亮即可。文件约25MB/张，建议有需求时再开。
  \item \textbf{安卓Motion Raw / Ultra RAW}：各品牌叫法不同。小米/OPPO/vivo旗舰机型在相机"专业模式"或"更多"里能找到RAW开关；华为在"专业模式"中开启RAW；三星Galaxy S系列在"专业"模式里有RAW选项。文件大小通常在20-50MB。
  \item \textbf{通用建议}：日常随手拍没必要开RAW，文件太大占空间；只有在重要场景（风光、人像、需要后期）时才开启。
\end{itemize}

\booksubsection{六、计算摄影原理}

现代手机摄影已经离不开“计算摄影”——也就是靠算法而不仅是镜头和传感器来生成照片。理解这几个核心原理，能让你更明白手机在按下快门的瞬间做了什么。

\subsubsection*{多帧合成：}

手机按下一次快门，其实拍了多张照片。算法把多帧对齐、加权融合，得到一张高质量成片。

\begin{itemize}
  \item \textbf{夜景模式}：1-3秒内连拍多张欠曝照片，对齐后叠加——噪点互相抵消、亮度累加。结果就是暗光也能拍出干净的画面。
  \item \textbf{HDR（高动态范围）}：分别拍摄欠曝、正常、过曝三张（或多张），合成一张亮部和暗部都有细节的照片。逆光场景尤其有效。
  \item \textbf{Deep Fusion（苹果）/Photonic Engine（光子引擎）}：针对中低光场景，逐像素分析多帧并融合，重点提升纹理和细节，毛衣、毛发这类细节会明显更清晰。
\end{itemize}

\subsubsection*{语义分割：}

AI不只是把整张图一刀切处理，而是先“看懂”画面里有什么——天空、人脸、草地、建筑、衣服——然后对每个区域分别应用不同的处理策略。

\begin{itemize}
  \item \textbf{天空}：单独增强蓝色、压低亮度，避免过曝。
  \item \textbf{人脸}：识别肤色区域，单独提亮、平滑肤色，不动五官和头发细节。
  \item \textbf{草地/树叶}：增强绿色饱和度，提升细节锐化。
  \item \textbf{背景虚化}：人像模式下，AI识别主体轮廓，背景模糊、主体清晰。
\end{itemize}

\subsubsection*{深度学习的应用：}

\begin{itemize}
  \item \textbf{AI场景识别}：相机自动识别"美食""宠物""夜景""文档"等场景，并对应优化参数。比如美食模式会增强暖色调和饱和度，让食物看起来更诱人。
  \item \textbf{AI人像虚化}：手机镜头物理上无法做到单反那种浅景深，靠深度学习模型预测主体边缘，再对背景做模糊。最新算法已经能正确处理头发丝、玻璃边缘这些细节。
  \item \textbf{AI消除/AI扩图}：选中画面中的路人或杂物，AI基于上下文"补全"被遮挡的部分。这是生成式AI在影像领域的应用。
\end{itemize}

\subsubsection*{计算摄影 vs 光学摄影：}

\begin{itemize}
  \item \textbf{计算摄影的优势}：小体积也能得到好画质，暗光、逆光等极端场景也能出片，普通用户不需要懂摄影知识就能拍出"好看"的照片。
  \item \textbf{计算摄影的劣势}：画面有时"过度处理"——锐化过度、色彩失真、天空颜色不自然；多帧合成在主体移动时会留下重影；细节是"算出来的"而非真实记录的。
  \item \textbf{光学摄影的优势}：大底传感器+大光圈镜头能真实捕捉光线，画质自然，后期空间大；不会因为算法"脑补"产生伪影。
  \item \textbf{趋势}：现代旗舰手机是"光学+计算"结合——用大底传感器获取真实信息，用算法优化呈现。两者并不对立。
\end{itemize}

\booksubsection{七、镜头类型与焦段}

现代旗舰手机通常配三到四颗镜头，每颗对应不同焦段。理解焦段，能让你在不同场景下选对镜头。

\subsubsection*{常见镜头与等效焦距：}

\begin{itemize}
  \item \textbf{超广角（13-24mm等效）}：视野最广，适合风光、建筑、室内空间。能拍到眼前所见的大场面，但边缘会有畸变（直线变弯），拍人像时人脸会被拉伸。适合"大场景"，不适合"特写"。
  \item \textbf{主摄（24-28mm等效）}：日常最常用的镜头，画质最好、光圈最大。兼顾视野和画质，拍人拍景都能用。大多数手机的"1x"就是主摄。
  \item \textbf{长焦（70-200mm等效）}：用于拍远景、人像特写。长焦能压缩空间感（远处的山和近处的人看起来更近），拍人像能获得更自然的五官比例和背景虚化。
  \item \textbf{潜望式长焦（200mm+等效）}：通过反射式光路设计，在手机厚度内实现更长焦距。能在很远距离拍清楚细节，适合演唱会、体育赛事、野生动物。旗舰机型才有。
\end{itemize}

\subsubsection*{传感器尺寸与像素大小的权衡：}

传感器是“底”，决定能接收多少光线。行业里说“底大一级压死人”——同样像素数下，传感器越大，单个像素也越大，进光量越多，画质越好。

\begin{itemize}
  \item \textbf{大底传感器}：进光量大，夜景画质好，动态范围宽。缺点是手机要做厚（镜头模组凸起），成本高。
  \item \textbf{高像素小底}：白天光线充足时，高像素能拍出更清晰的照片，方便后期裁切。但暗光下每个像素进光不足，噪点多。
  \item \textbf{主流旗舰选择}：主摄多为1/1.3英寸到1英寸（部分机型已用上1英寸大底），像素数5000万左右——这是画质和体积的平衡点。
\end{itemize}

\subsubsection*{像素融合技术（Pixel Binning）：}

为了兼顾白天高像素和夜晚大像素，手机厂商用“像素融合”——把多个小像素合并成一个大像素输出。

\begin{itemize}
  \item \textbf{四合一像素（Quad Bayer）}：5000万像素传感器，输出1250万像素照片。每个输出像素由4个相邻像素合成，进光量提升4倍，画质更干净。
  \item \textbf{九合一像素（Nonacell）}：用于1亿像素以上的传感器，输出约1100-1200万像素。同样原理，进光量更大。
  \item \textbf{实际效果}：默认拍照模式下，你拿到的就是融合后的低像素照片——画质好、文件小。需要极致细节时（比如拍文档、远景），可以手动切到"高像素模式"输出全分辨率。
\end{itemize}

\booksubsection{八、特殊拍摄模式}

除了日常拍照，手机里还有几个特殊模式，用对了能拍到平时拍不出的画面。

\subsubsection*{全景模式：}

\begin{itemize}
  \item 拍摄时按提示箭头缓慢平稳地横向移动手机，不要快——快了会出现拼接错位。
  \item 保持手机垂直方向稳定，只横向平移。如果上下晃动，地平线会变波浪。
  \item 有移动的人或车时慎用——拼接时会出现"半个身体"或"鬼影"。
  \item 全景模式也可以竖着拍——竖拍横拼接，适合拍高耸的建筑或瀑布。
\end{itemize}

\subsubsection*{慢动作与延时摄影：}

\begin{itemize}
  \item \textbf{慢动作}：用120fps或240fps高帧率拍摄，回放时变成慢镜头。适合拍水滴、运动瞬间、宠物跳跃。光线要充足，慢动作模式下进光时间短，暗光画质差。
  \item \textbf{延时摄影}：间隔拍摄多张照片合成快进视频。适合拍车流、云彩移动、日落、花朵开放。拍摄时长几分钟到几小时不等，务必用三脚架固定手机，手持必糊。
  \item \textbf{小技巧}：延时摄影拍云彩时，把手机架好，拍15-30分钟，最后得到10秒左右的延时视频，效果震撼。
\end{itemize}

\subsubsection*{星空/天文模式：}

\begin{itemize}
  \item 需要远离城市光污染，找完全没有路灯干扰的郊外或山区。城市里拍不出银河。
  \item 必须用三脚架，手持必糊。手机要稳稳架住几秒到几十秒。
  \item 天气要晴朗无云，月亮不要太大（满月会盖住星光）。农历月初或月末的暗夜最适合。
  \item iPhone的"夜间模式"会自动识别星空并延长曝光时间（最长30秒）；安卓旗舰的"星空模式"或"天文模式"效果类似。要耐心等待，按下快门后不要动手机。
\end{itemize}

\subsubsection*{微距模式：}

\begin{itemize}
  \item 用于拍摄花朵、昆虫、珠宝、文字特写等小物体。能贴近到1-2厘米距离对焦。
  \item 部分手机有专门的微距镜头（多为200万像素或500万像素的小底镜头），画质一般；旗舰机用超广角镜头兼顾微距，画质更好。
  \item 贴近拍摄时镜头离主体很近，会遮挡光线，必要时用台灯或手电补光。
  \item 对焦距离极近，景深很浅，主体前后移动几毫米就可能脱焦。拍摄时手要稳，最好让主体保持不动。
\end{itemize}

\begin{tcbinfo}{静态摄影实用建议}
\begin{itemize}
  \item 日常拍摄：HEIF/JPEG默认即可，文件小、画质够用。
  \item 重要场景（风光、人像、需要后期）：开启RAW，保留最大后期空间。
  \item 选对镜头：风光用超广角，人像用长焦，日常用主摄。
  \item 夜景和星空：务必固定手机，开夜景模式或天文模式。
  \item 大底手机（如1英寸传感器机型）画质上限更高，但操作习惯仍是关键。
\end{itemize}

\end{tcbinfo}

\begin{tcbpractice}{动手试试：摸清你的相机能力}
\begin{enumerate}
  \item 查看你的手机是否支持RAW拍摄（设置相机格式），试着开启并拍摄一张，对比同一场景的JPEG和RAW文件大小。
  \item 切换到超广角、主摄、长焦三个焦段，分别拍同一处远景，体会视角差异和画质差异。
  \item 找一个有夜景模式或星空模式的场景，固定手机后拍摄一张夜景照片。
  \item 尝试用全景模式横拍一张大场景，注意保持手机平稳移动。
\end{enumerate}

\end{tcbpractice}
